% Original analytic teaching diagram, 2026-11-12.
% Huffman tree for the indices of the 4-level Lloyd-Max quantizer of Laplace(0,1):
% p = (0.1016, 0.3984, 0.3984, 0.1016) for j = 1..4. Merges: 1+4 -> 0.2032; 0.2032+2 -> 0.6016; +3 -> 1.
% Codewords: j=3: 0, j=2: 10, j=1: 110, j=4: 111. Average length 1.8048 bit.
\definecolor{VB}{HTML}{7188A3}\definecolor{VBF}{HTML}{E5EAF0}
\definecolor{VG}{HTML}{829C8B}\definecolor{VGF}{HTML}{E8EEE9}
\definecolor{VY}{HTML}{B4A16F}\definecolor{VYF}{HTML}{F1ECDD}
\definecolor{VR}{HTML}{A7787D}\definecolor{VRF}{HTML}{F0E5E6}
\definecolor{VP}{HTML}{9588A5}\definecolor{VPF}{HTML}{ECE8F0}
\definecolor{VI}{HTML}{3B4252}\definecolor{VM}{HTML}{6B7280}\definecolor{VL}{HTML}{4C4D4F}
\begin{tikzpicture}[x=1mm,y=1mm,>=Stealth,font=\figurefont\fontsize{8}{10}\selectfont,
lab/.style={text=VM,font=\figurefont\fontsize{7.5}{9}\selectfont},
nd/.style={circle,draw=VL,fill=white,line width=.6pt,inner sep=1.2pt,text=VI},
lf/.style={rectangle,draw=VB,fill=VBF,line width=.6pt,minimum width=13mm,minimum height=6mm,text=VI},
e/.style={draw=VL,line width=.6pt}]
\node[nd] (r) at (0,0) {$1$};
\node[lf] (j3) at (-18,-12) {$j=3$};
\node[nd] (a) at (18,-12) {$0.6016$};
\node[lf] (j2) at (2,-25) {$j=2$};
\node[nd] (b) at (34,-25) {$0.2032$};
\node[lf] (j1) at (20,-38) {$j=1$};
\node[lf] (j4) at (48,-38) {$j=4$};
\draw[e] (r) -- node[above left,text=VG] {$0$} (j3);
\draw[e] (r) -- node[above right,text=VG] {$1$} (a);
\draw[e] (a) -- node[above left,text=VG] {$0$} (j2);
\draw[e] (a) -- node[above right,text=VG] {$1$} (b);
\draw[e] (b) -- node[above left,text=VG] {$0$} (j1);
\draw[e] (b) -- node[above right,text=VG] {$1$} (j4);
\node[lab] at (-18,-18) {$0.3984$\quad\texttt{0}};
\node[lab] at (2,-31) {$0.3984$\quad\texttt{10}};
\node[lab] at (20,-44) {$0.1016$\quad\texttt{110}};
\node[lab] at (48,-44) {$0.1016$\quad\texttt{111}};
\end{tikzpicture}
